%%=============================================================================
%% Voorwoord
%%=============================================================================

\chapter*{\IfLanguageName{dutch}{Woord vooraf}{Preface}}%
\label{ch:voorwoord}

Deze graduaatsproef is het resultaat van mijn stage bij Be-Mobile NV, een Belgisch bedrijf gespecialiseerd in verkeersdata en navigatieoplossingen voor partners zoals Flitsmeister en TomTom. Tijdens mijn stage werkte ik in het platformteam, waar ik de opdracht kreeg om observability op te zetten voor hun nieuwe API gateway op basis van APISIX.

De leercurve was steil: OpenTelemetry, Prometheus, Thanos, Loki, Tempo en Grafana Alloy --- elk met hun eigen configuratielogica en query-taal. Ik hoop dat dit verslag niet alleen de technische resultaten documenteert, maar ook de redenering achter de keuzes helder maakt voor iedereen die een gelijkaardige integratie wil uitvoeren.

Ik wil graag Glenn Delanghe bij Be-Mobile bedanken voor de vrijheid en het vertrouwen die ik kreeg om dit onderzoek zelfstandig uit te voeren, en voor de waardevolle feedback tijdens de wekelijkse check-ins. Ook dank aan Dirk Vandycke aan HOGENT voor de begeleiding bij het uitschrijven van het onderzoeksvoorstel en het opvolgen van de voortgang.

De broncode, configuratiebestanden en Grafana-dashboards die bij dit onderzoek horen, zijn publiek beschikbaar via de GitHub-repository:
\url{https://github.com/jipgg/graduaat-programmeren-graduaatsproef}